% NOTE: mnras.cls not available locally — using article for draft compilation.
% Switch to \documentclass[fleqn,usenatbib]{mnras} for journal submission.
\documentclass[12pt,a4paper]{article}

\usepackage{amsmath}
\usepackage{graphicx}
\usepackage{booktabs}
\usepackage{hyperref}
\usepackage{natbib}
\usepackage{geometry}
\geometry{margin=1in}

\title{The Hubble Tension as a Measurement Artifact II:\\
Residual Decomposition of $\xi$-Normalized BAO Constraints\\
and the $\Omega_m$ Deficit}

\author{Eric D. Martin\\
\small Independent Researcher}

\date{2026-03-25\\
\small Preprint --- Session: bdf10532-f890-40c7-8abf-adaa29ec06b1}

\begin{document}
\maketitle

\begin{abstract}
The reported $5\sigma$ Hubble tension is widely interpreted as a breakdown of
the $\Lambda$CDM model. By applying Horizon-Normalized Coordinates
($\xi = d_c/d_H$), we demonstrate that this tension decomposes into two
distinct, separable layers rather than a single monolithic discrepancy. First,
we confirm that $\sim$90\% of the apparent tension is a ``frame-mixing''
artifact---a mathematical byproduct of comparing datasets in incompatible
$H_0$-dependent coordinate systems \citep{Martin2025_Paper1}. Once this
artifact is removed via $\xi$-normalization, a physical residual remains.
Using the DESI 2024/25 Baryon Acoustic Oscillation (BAO) data
\citep{DESI_2024_DR1, DESI_2025_DR2}, we fit a $\chi^2$ minimizer with
self-consistent sound-horizon scaling $r_s(\Omega_m)$ and find a best-fit
$\Omega_m = 0.268$ (7-bin $D_M/r_s$ analysis, $\Delta\chi^2 = 7.06$ vs.\
Planck $\Lambda$CDM), confirming the direction of the matter density deficit
independently reported by DESI DR1 ($\Omega_m = 0.295 \pm 0.010$). Adding
$w_0w_a$ freedom reduces $\chi^2$ by only $0.12$ for two additional parameters,
providing no statistical support for dynamical dark energy in this dataset
alone. We conclude that the Hubble tension is largely a measurement artifact;
its removal reveals a genuine matter density deficit consistent across
independent probes, but does not require evolving dark energy from current BAO
data. Confirmation of both signals awaits Euclid DR1.
\end{abstract}

% ─────────────────────────────────────────────────────────────────────────────
\section{Introduction}

The Hubble tension has persisted for over a decade, with the Planck CMB
analysis yielding $H_0 = 67.4 \pm 0.5$ km/s/Mpc \citep{Planck2020} and the
SH0ES distance ladder yielding $H_0 = 73.04 \pm 1.04$ km/s/Mpc
\citep{Riess2022}. The $5\sigma$ discrepancy has motivated proposals for early
dark energy \citep{Poulin2019}, modified gravity, and non-standard
recombination. In \citet{Martin2025_Paper1}, we demonstrated that 93.3\% of
this tension is a frame-mixing artifact: a mathematical consequence of
comparing datasets in incompatible $H_0$-dependent coordinate systems. The
residual physical signal is a $\sim$3\% matter density discrepancy.

Here we extend that analysis to characterise the residual using the DESI
2024/25 BAO dataset \citep{DESI_2024_DR1, DESI_2024_w0wa, DESI_2025_DR2}. We
find that the residual is well described by a global $\Omega_m$ deficit,
confirmed across multiple independent probes. Within current BAO data, adding
$w_0w_a$ parameters does not improve the fit significantly ($\Delta\chi^2 <
0.2$), leaving the dark energy question open for future datasets.

% ─────────────────────────────────────────────────────────────────────────────
\section{Method: $\xi$-Normalization Applied to BAO}

We apply the horizon-normalized coordinate framework from
\citet{Martin2025_Paper1}:
\begin{equation}
    \xi(z) = \frac{d_c(z)}{d_H} = \int_0^z \frac{dz'}{E(z')}
\end{equation}
where for the general $w_0w_a$ dark energy parametrization (CPL;
\citealt{Chevallier2001, Linder2003}):
\begin{equation}
    E(z) = \sqrt{\Omega_m(1+z)^3 + (1-\Omega_m)(1+z)^{3(1+w_0+w_a)}
    e^{-3w_a z/(1+z)}}
\end{equation}
For redshift-derived distances, $d_c \propto 1/H_0$ and $d_H \propto 1/H_0$,
so $\xi$ is exactly $H_0$-independent. We apply this to all 7 DESI DR1 BAO
bins spanning $0.295 \leq z \leq 2.33$ \citep{DESI_2024_DR1}.

A critical methodological requirement is that the sound horizon $r_s$ must be
varied self-consistently with $\Omega_m$. Fixing $r_s$ at the Planck fiducial
(147.09 Mpc) while varying $\Omega_m$ inverts the sign of the fit, producing a
spurious apparent excess rather than the physical deficit. We use the
Percival~(2007) scaling:
\begin{equation}
    r_s(\Omega_m) = 147.09 \left(\frac{\Omega_m h^2}{0.1430}\right)^{-0.255}
    \text{ Mpc}
\end{equation}
normalized to the Planck~2018 value at $\Omega_m = 0.315$, $h = 0.674$
\citep{Planck2020}. This ensures $D_M/r_s$ is computed with both numerator and
denominator varying as functions of $\Omega_m$, preserving the $H_0$-independence
of the ratio.

We run $\chi^2$ minimizers over $\{\Omega_m\}$, $\{\Omega_m, w_a\}$, and
$\{\Omega_m, w_0, w_a\}$ to quantify each parameter's contribution.

% ─────────────────────────────────────────────────────────────────────────────
\section{Results}

\subsection{$\Omega_m$ Deficit: Global Shift}

Applying $\xi$-normalization to all 7 DESI DR1 bins with self-consistent
$r_s(\Omega_m)$ scaling, we find a systematic offset between the observed
$D_M/r_s$ values and the Planck $\Lambda$CDM prediction. A 1D $\chi^2$
minimizer yields a best-fit $\Omega_m = 0.268$ with $\chi^2 = 6.51$ vs.\ 13.57
for Planck, a $\Delta\chi^2 = 7.06$ improvement. The matter density deficit
($\Omega_m < 0.315$) is confirmed in the same direction by DESI DR1 BAO alone
($\Omega_m = 0.295 \pm 0.010$; \citealt{DESI_2024_DR1}), the DESI 2026
combined analysis ($\Omega_m \approx 0.304$; \citealt{DESI_2025_DR2}), and the
$S_8/\sigma_8$ weak lensing signal from KiDS-1000 and DES~Y3
\citep{Heymans2021, Amon2022}.

Our best-fit $\Omega_m = 0.268$ is lower than the DESI DR1 full-analysis value
of $0.295 \pm 0.010$. This discrepancy likely reflects the limitations of the
Percival (2007) $r_s$ approximation and our restriction to $D_M/r_s$ data
without the full covariance matrix (see Section~\ref{sec:limitations}). The
qualitative conclusion is robust: the matter density is below the Planck
$\Lambda$CDM value across all analyses.

\subsection{The $z = 0.51$ Pivot and the Absence of a $w_0w_a$ Signal}
\label{sec:z051}

The LRG1 bin at $z = 0.51$ shows a persistent residual of $+1.80\sigma$ at
the corrected best-fit $\Omega_m = 0.268$. This single-bin excess could reflect
systematic effects in the LRG1 photometric calibration, an unmodelled
selection-function feature, or a genuine physical anomaly. However, we
emphasise that it does not constitute evidence for dynamical dark energy within
the current dataset.

Adding the full $w_0w_a$ parametrization to the fit reduces $\chi^2$ by only
$0.12$ relative to the 1D $\Omega_m$-only fit (Table~\ref{tab:chi2}). For two
additional parameters, this corresponds to $p = 0.94$: the $w_0w_a$ degrees of
freedom are entirely unsupported by the data. The apparent LRG1 anomaly is
absorbed into the $\Omega_m$ correction; no oscillatory multi-bin signature
indicative of a CPL equation-of-state bend is present in the residuals.

\begin{table}[h]
\centering
\caption{$\chi^2$ goodness-of-fit comparison across 7 DESI DR1 BAO bins.
All $\chi^2$ values from \texttt{bao\_wa\_minimizer.py} with
self-consistent $r_s(\Omega_m)$ scaling.}
\label{tab:chi2}
\begin{tabular}{lcccccc}
\toprule
Model & $\Omega_m$ & $w_0$ & $w_a$ & $\chi^2$ & $\chi^2_\nu$ &
Interpretation \\
\midrule
Planck $\Lambda$CDM        & 0.315 & $-1.0$ & $\phantom{+}0.0$   & 13.57 & 2.26 &
  Rejected \\
$\Omega_m$ free (UHA)     & 0.268 & $-1.0$ & $\phantom{+}0.0$   & \phantom{0}6.51 & 1.09 &
  Good fit \\
$\Omega_m + w_a$          & 0.264 & $-1.0$ & $-0.081$ & \phantom{0}6.42 & 1.28 &
  No improvement \\
$\Omega_m + w_0 + w_a$    & 0.251 & $-1.101$ & $+0.360$ & \phantom{0}6.39 & 1.60 &
  No improvement \\
\bottomrule
\end{tabular}
\end{table}

% ─────────────────────────────────────────────────────────────────────────────
\section{Discussion}

\subsection{The Double Hierarchy}

The Hubble tension decomposes into two separable layers under
$\xi$-normalization:
\begin{enumerate}
    \item \textbf{$\sim$90\% coordinate artifact} --- removed by
    $\xi$-normalization; the $H_0$ factor cancels exactly for redshift-derived
    distances.
    \item \textbf{$\sim$5\% matter density deficit} --- a global
    $\Omega_m < 0.315$ confirmed independently by $\xi$-residuals (this work),
    DESI DR1 BAO \citep{DESI_2024_DR1}, DESI DR2 \citep{DESI_2025_DR2}, and
    $S_8/\sigma_8$ weak lensing \citep{Heymans2021, Amon2022}.
\end{enumerate}

A third layer --- evolving dark energy ($w_0w_a$) --- was suggested by an
earlier, incorrect analysis in which $r_s$ was held fixed at the Planck
fiducial while $\Omega_m$ was varied. Fixing $r_s$ inverts the direction of
the $\Omega_m$ sensitivity and produces a spurious $w_0w_a$ improvement
($\Delta\chi^2 = 4.74$ in the incorrect version). With the corrected
self-consistent $r_s(\Omega_m)$ scaling, the $w_0w_a$ improvement collapses to
$\Delta\chi^2 = 0.12$. We report the corrected result here; the dark energy
question remains open and requires Euclid DR1 to resolve.

\subsection{The $S_8$ Tension as a Matter Density Tie-Breaker}

The most compelling independent evidence for a physical $\Omega_m$ deficit will
arrive with the Euclid DR1 $S_8$ results \citep{Euclid2022_Survey}. The $S_8$
parameter, $\sigma_8(\Omega_m/0.3)^{0.5}$, measures the amplitude of matter
clustering and has consistently trended $2$--$3\sigma$ below the Planck
$\Lambda$CDM prediction in prior weak lensing surveys \citep{Heymans2021,
Amon2022}. Within our $\xi$-normalized framework, this clustering tension is
not a separate crisis but a secondary manifestation of the same $\Omega_m$
deficit. A genuine deficit in $\Omega_m$ relative to Planck mathematically
forces a lower $S_8$ value. Confirmation of $S_8 < 0.83$ at high significance
would constitute a definitive tie-breaker for the matter density hypothesis
over a coordinate artifact.

\subsection{Reply to the JWST $8\sigma$ Objection}

The JWST NIRCam result \citep{Riess2023} confirming $H_0 = 73.17 \pm 1.01$
km/s/Mpc is sometimes cited as evidence that the tension cannot be a coordinate
artifact. As shown in \citet{Martin2025_Paper1}, JWST measures the artifact
more precisely within a fixed $H_0$ frame; it does not test whether that frame
and the Planck frame occupy compatible coordinates. The $\Omega_m$ deficit
identified in this work emerges from the internal geometry of the DESI BAO
dataset, which is $H_0$-independent by construction and therefore orthogonal to
the systematic uncertainties of the local distance ladder.

\subsection{Limitations and Non-Claims}
\label{sec:limitations}

This analysis has several important limitations. (1)~The $\chi^2$ minimizer
uses only the $D_M/r_s$ observables; a full analysis would incorporate
$D_H/r_s$ and the complete bin covariance matrix, which is expected to shift
the $\Omega_m$ central value toward the DESI-reported 0.295. (2)~The
$r_s(\Omega_m)$ scaling from Percival~(2007) is an approximation; a more
accurate treatment would vary both $\Omega_b h^2$ and $\Omega_m h^2$ jointly.
(3)~We make no claim of a detection of dynamical dark energy from this
analysis; the $\Delta\chi^2 = 0.12$ for $w_0w_a$ is statistically negligible.
(4)~The LRG1 $z=0.51$ residual ($+1.80\sigma$) is unexplained but insufficient
to establish CPL evolution. (5)~All conclusions are conditioned on the DESI DR1
public data release; updated constraints from DR2 and independent datasets may
shift quantitative results.

% ─────────────────────────────────────────────────────────────────────────────
\section{Conclusion}

The Hubble tension decomposes into two separable layers under $\xi$-normalized
BAO analysis. The dominant layer ($\sim$90\%) is a coordinate artifact removed
by horizon-normalization. The residual is a matter density deficit confirmed
independently by DESI BAO, weak lensing surveys, and the $\xi$-normalized
analysis presented here. Current BAO data do not support a dynamical dark
energy component; the $w_0w_a$ improvement is statistically negligible
($\Delta\chi^2 < 0.2$ for two parameters). The $\xi$-normalization framework
provides the infrastructure to separate coordinate noise from physical signals
as higher-precision datasets arrive.

% ─────────────────────────────────────────────────────────────────────────────
\section*{Competing Interests}

The author holds US Provisional Patent 63/902,536 (Universal Horizon Address,
UHA) covering the $\xi$ coordinate infrastructure described in this work
(Zenodo DOI:~10.5281/zenodo.17435574). The theoretical results reported here
are independent of the patent claims and do not constitute commercial promotion.

% ─────────────────────────────────────────────────────────────────────────────
\section*{Data Availability}

All analysis code is publicly available at \url{https://github.com/abba-01/uha}.
Reproduction scripts: \texttt{bao\_wa\_discriminator.py},
\texttt{bao\_wa\_minimizer.py}, \texttt{uha\_evolution\_plot.py}.
The DESI DR1 dataset is the public release from \citet{DESI_2024_DR1}.
This manuscript is deposited at Zenodo: DOI~10.5281/zenodo.19272498.
No proprietary data were used.

% ─────────────────────────────────────────────────────────────────────────────
\bibliographystyle{mnras}
\bibliography{paper2_refs}

% Manual fallback bibliography (if mnras.bst unavailable at initial submission)
\begin{thebibliography}{99}

\bibitem[Amon et al.(2022)]{Amon2022}
Amon A., et al., 2022, Phys.\ Rev.\ D, 105, 023514,
DOI: 10.1103/PhysRevD.105.023514

\bibitem[Brout et al.(2022)]{Brout2022}
Brout D., et al., 2022, ApJ, 938, 110,
DOI: 10.3847/1538-4357/ac8e04

\bibitem[Chevallier \& Polarski(2001)]{Chevallier2001}
Chevallier M., Polarski D., 2001, Int.\ J.\ Mod.\ Phys.\ D, 10, 213,
DOI: 10.1142/S0218271801000822

\bibitem[DESI Collaboration(2024a)]{DESI_2024_DR1}
DESI Collaboration, 2024a, arXiv:2404.08633

\bibitem[DESI Collaboration(2024b)]{DESI_2024_w0wa}
DESI Collaboration, 2024b, arXiv:2404.08056

\bibitem[DESI Collaboration(2025)]{DESI_2025_DR2}
DESI Collaboration, 2025, arXiv:2503.10806

\bibitem[Euclid Collaboration et al.(2022)]{Euclid2022_Survey}
Euclid Collaboration, Scaramella R., et al., 2022, A\&A, 662, A112,
DOI: 10.1051/0004-6361/202141938

\bibitem[Freedman(2024)]{Freedman2024}
Freedman W.~L., 2024, ApJ, 961, 32,
DOI: 10.3847/1538-4357/ad1bac

\bibitem[Heymans et al.(2021)]{Heymans2021}
Heymans C., et al., 2021, A\&A, 646, A140,
DOI: 10.1051/0004-6361/202014214

\bibitem[Linder(2003)]{Linder2003}
Linder E.~V., 2003, Phys.\ Rev.\ Lett., 90, 091301,
DOI: 10.1103/PhysRevLett.90.091301

\bibitem[Martin(2025a)]{Martin2025_UHA}
Martin E.~D., 2025a, Universal Horizon Address, US Patent 63/902,536,
DOI: 10.5281/zenodo.17435574

\bibitem[Martin(2025b)]{Martin2025_Paper1}
Martin E.~D., 2025b, The Hubble Tension as a Measurement Artifact,
DOI: 10.5281/zenodo.19211662

\bibitem[Planck Collaboration(2020)]{Planck2020}
Planck Collaboration, 2020, A\&A, 641, A6,
DOI: 10.1051/0004-6361/201833910

\bibitem[Poulin et al.(2019)]{Poulin2019}
Poulin V., et al., 2019, Phys.\ Rev.\ Lett., 122, 221301,
DOI: 10.1103/PhysRevLett.122.221301

\bibitem[Riess et al.(2022)]{Riess2022}
Riess A.~G., et al., 2022, ApJL, 934, L7,
DOI: 10.3847/2041-8213/ac5c5b

\bibitem[Riess et al.(2023)]{Riess2023}
Riess A.~G., et al., 2023, ApJL, 956, L18,
DOI: 10.3847/2041-8213/acf769

\end{thebibliography}

\end{document}
